The Lake (LAK) Package computes the stage of each lake from a water balance: precipitation, evaporation, runoff, specified inflow and withdrawal, outlet flow, lake storage, and the lakebed seepage exchanged with the connected groundwater flow (GWF) cells \citep[\S 8]{modflow6gwf}. By default the package solves this balance for the stage with a substitution iteration nested inside each outer (Picard or Newton) iteration of the GWF solution: the stage is updated to balance its budget against the latest GWF heads, and the resulting seepage is then added to the GWF matrix. Because the stage lags the heads by one outer iteration, strongly coupled and steady-state lakes can require many outer iterations to converge, and some do not converge at all.

The \texttt{IMPLICIT} option instead adds each lake stage to the GWF matrix as an extra unknown (one row and column) and solves it together with the GWF heads. The lakebed seepage to each connected cell is assembled as a conductance coupling between the lake-stage row and the cell row. In the seepage, the lake stage and the connected-cell head are each kept greater than or equal to the lake bottom, smoothed with a ramp function near the bottom, so the seepage and its derivatives change continuously as a connection wets or dries. The remaining (non-seepage) budget terms are linearized in stage with a finite difference and added to the diagonal and right-hand side. If Newton under-relaxation is active, a stage driven below the lake bottom during an iteration is relaxed back toward the bottom. Because the stage is no longer lagged, the coupled system typically converges in far fewer outer iterations -- often one to three orders of magnitude fewer for steady-state and strongly coupled lakes -- and problems that fail to converge with the default formulation are routinely solved. At convergence the two formulations satisfy the same water balance and give equivalent heads, stages, and budgets.

\subsection{Perched (Disconnected) Lakes}

A lake is perched when the groundwater head at every connected cell falls below the lakebed, so all of its seepage is downward leakage. The lake-stage equation stays well posed only if its diagonal -- the derivative of the lake balance with respect to stage -- remains non-zero. Holding the connected-cell head at the lake bottom, once it has fallen below the lakebed, keeps this leakage dependent on the lake stage, because a deeper lake drives more leakage through its bed. The diagonal therefore stays non-zero, and the implicit formulation determines the stage even when the lake is completely disconnected from the aquifer in a steady-state stress period. The stage it finds matches the default formulation, which balances the same leakage.

One case still favors the default formulation: a lake whose connected cells dry and rewet under the standard rewetting (REWET) capability rather than the Newton formulation. Rewetting can settle on more than one valid configuration, so the implicit and default approaches, while both mass conserving, may reach different results. The Newton formulation keeps cells active and avoids this ambiguity, and is recommended for implicit lakes that may become perched.
